%   Il progetto nasce dal template per il frontespizio realizzato da Marco Antonio Corallo, che ringrazio. Seguono alcuni commenti per evidenziare la presenza di alcuni pacchetti che mi sono stati utili per la stesura della tesi. Chiaramente, dipende tutto dal tipo di lavoro che uno vuole eseguire, che determina anche le diverse esigenze. Durante la stesura ho passato molto tempo su siti e forum a cercare di risolvere alcuni probelmi di formattazione, ma in generale Latex è stato piuttosto versatile. 

% Tipo di documento. L'uso di twoside implica che i capitoli inizino sempre con la prima pagina a sinistra, eventualmente lasciando una pagina vuota nel capitolo precedente. Se questa cosa è fastidiosa, è possibile rimuoverlo. 
\documentclass[a4paper, twoside,openright]{report}
\usepackage{siunitx}

\DeclareMathSizes{10}{10}{7}{4}

% Dimensione dei margini
\usepackage[a4paper,top=3cm,bottom=3cm,left=3cm,right=3cm]{geometry} 
% Dimensione del font
\usepackage[fontsize=13pt]{scrextend}
% Lingua del testo
\usepackage[english,italian]{babel}

\usepackage{natbib}
% \bibpunct{(}{)}{;}{a}{}{,} % to follow the A&A style

% Lingua per la bibliografia
\usepackage[fixlanguage]{babelbib}
% Codifica del testo
\usepackage[utf8]{inputenc} 
% Encoding del testo
\usepackage[T1]{fontenc}

\usepackage[italian]{babel} % qui si dichiara la lingua del documento

% Permette di generare testo fittizio. Mi è stato utile 
% per capire quale sarebbe stata l'impostazione del 
% testo nella pagina prima che scrivessi un determinato paragrafo
\usepackage{lipsum}
% Per ruotare le immagini
\usepackage{rotating}
% Per modificare l'header delle pagine 
\usepackage{fancyhdr}    
\usepackage{booktabs}
\usepackage{float}

% Librerie matematiche
\usepackage{amssymb}
\usepackage{amsmath}
\usepackage{amsthm} 
\usepackage{accents}
\usepackage{aas_macros}

% Uso delle immagini
\usepackage{graphicx}
% Uso dei colori
\usepackage[dvipsnames]{xcolor}         
% Uso dei listing per il codice
\usepackage{listings}          
% Per inserire gli hyperlinks tra i vari elementi del testo 
\usepackage{hyperref}     
% Diversi tipi di sottolineature
\usepackage[normalem]{ulem}

\usepackage{xspace}
% -----------------------------------------------------------------

% Modifica lo stile dell'header
\pagestyle{fancy}
\fancyhf{}
\lhead{\rightmark}
\rhead{\textbf{\thepage}}
\fancyfoot{}
\setlength{\headheight}{12.5pt}

% Rimuove il numero di pagina all'inizio dei capitoli
\fancypagestyle{plain}{
  \fancyfoot{}
  \fancyhead{}
  \renewcommand{\headrulewidth}{0pt}
}

% Stile del codice
\lstdefinestyle{codeStyle}{
    % Colore dei commenti
    commentstyle=\color{black},
    % Colore delle keyword
    keywordstyle=\color{black},
    % Stile dei numeri di riga
    numberstyle=\tiny\color{black},
    % Colore delle stringhe
    stringstyle=\color{black},
    % Dimensione e stile del testo
    basicstyle=\ttfamily\footnotesize,
    % newline solo ai whitespaces
    breakatwhitespace=false,     
    % newline si/no
    breaklines=true,                 
    % Posizione della caption, top/bottom 
    captionpos=b,                    
    % Mantiene gli spazi nel codice, utile per l'indentazione
    keepspaces=true,                 
    % Dove visualizzare i numeri di linea
    numbers=left,                    
    % Distanza tra i numeri di linea
    numbersep=5pt,                  
    % Mostra gli spazi bianchi o meno
    showspaces=false,                
    % Mostra gli spazi bianchi nelle stringhe
    showstringspaces=false,
    % Mostra i tab
    showtabs=false,
    % Dimensione dei tab
    tabsize=2
} \lstset{style=codeStyle}

% Stile di codice per dimensioni maggiori, in cui ho avuto bisogno di un testo più picolo (ad esempio se si vuole inserire del codice che ha linee molto lunghe). Per usare questo stile piuttosto che il precedente, usare 

% \lstset{style=longBlock}
%  % inserire il codice...
% \lstset{style=codeStyle}

% Il secondo comando consente di tornare allo stile precedente 
\lstdefinestyle{longBlock}{
    commentstyle=\color{black},
    keywordstyle=\color{black},
    numberstyle=\tiny\color{black},
    stringstyle=\color{black},
    basicstyle=\ttfamily\scriptsize,
    breakatwhitespace=false,         
    breaklines=true,                 
    captionpos=b,                    
    keepspaces=true,                 
    numbers=left,                    
    numbersep=5pt,                  
    showspaces=false,                
    showstringspaces=false,
    showtabs=false,                  
    tabsize=2
} \lstset{style=codeStyle}

% Togliendo il commento al comando che segue, si inseriscono nella bibliografia anche le fonti presenti in Bibliography.bib ma non citati direttamente con il comando \cite
% \nocite{*}

% Margini prima e dopo blocchi di codice, per avere più distanza
\lstset{aboveskip=20pt,belowskip=20pt}

% Modifica dello stile dei riferimenti, con il testo in cyano
%\hypersetup{
%    colorlinks,
%    linkcolor=CornflowerBlue,
%    citecolor=CornflowerBlue
%}

% Aggiunti definizioni, teoremi, linea e listing
\newtheorem{definition}{Definizione}[section]
\newtheorem{theorem}{Teorema}[section]
\providecommand*\definitionautorefname{Definizione}
\providecommand*\theoremautorefname{Teorema}
\providecommand*{\listingautorefname}{Listing}
\providecommand*\lstnumberautorefname{Linea}

\raggedbottom

%\newcommand{\cgs}[1]{{\textcolor{brown}[\textcolor{red}{\bf{GS: }}{ \textcolor{brown}{#1]}}}}             
%\newcommand{\cmc}[1]{{\textcolor{blue}[\textcolor{magenta}{\bf{MC: }}{ \textcolor{blue}{#1]}}}}


% User defined commands
\newcommand{\src}{MAXI J1816--195\xspace}

% -----------------------------------------------------------------
\begin{document}

\begin{titlepage}
\begin{figure}[!htb]
    \centering
    \includegraphics[keepaspectratio=true,scale=0.5]{images/Frontespizio/download.png}
\end{figure}

\begin{center}
    \LARGE{Università degli Studi di Cagliari}
    \vspace{5mm}
    \\ \large{Facoltà di Scienze Matematiche, Fisiche e Naturali}
    \vspace{5mm}
    \\ \LARGE{Laurea Magistrale in Fisica}
\end{center}

\vspace{15mm}
\begin{center}
    {\LARGE{\bf Thesis Title}}
    
    % Se il titolo è abbastanza corto da stare su una riga, si può usare
    
    % {\LARGE{\bf Un fantastico titolo per la mia tesi!}}
\end{center}
\vspace{30mm}

\begin{minipage}[t]{0.47\textwidth}
	{\large{Relatore:}{\normalsize\vspace{3mm}
	\\ 
        \large{\textbf{Prof. Riccardo Murgia}}\vspace{3mm} \\
        {\large{Co-relatori:}{\normalsize\vspace{3mm}
 	\bf\\ 
        \large{Prof. Thomas Kupfer\vspace{1mm}\\
                     Weitian Yu}
    \normalsize\vspace{3mm}\bf}}}}\\ 
\end{minipage}
\hfill
\begin{minipage}[t]{0.47\textwidth}\raggedleft
	{\large{Candidato:}{\normalsize\vspace{3mm} \bf\\ \large{Gabriele Todde}}}
\end{minipage}

\vspace{30mm}
\hrulefill
\\\centering{\large{ANNO ACCADEMICO 2024/2025}}

\end{titlepage}

\begin{abstract}
\addcontentsline{toc}{chapter}{Abstract}
In recent years ground-based interferometers such as LIGO and Virgo opened the high-frequency window of the rapidly evolving gravitational wave astronomy.
Future space-based detectors, in particular the \textit{Laser Interferometer Space Antenna} (LISA), will extend gravitational wave detection into the low-frequency band, where Galactic compact binaries such as double white dwarfs (DWDs) are expected to shape LISA's sensitivity curve significantly;
the aim of this work is to investigate the expected gravitational wave background produced by \textit{extragalactic} DWDs in the local Universe, employing the \textit{Compact Object Synthesis and Monte Carlo Investigation Code} (\textbf{COSMIC}) to generate synthetic populations representing the galaxies included in the observational survey \textbf{Gravitational Wave Galaxy Catalog} (GWGC).
Each simulated binary contributes to the total gravitational wave strain, which is then combined across frequency bins and compared against the LISA sensitivity curve.
The results show that the total simulated extragalactic signal lies approximately one order of magnitude below LISA’s detection threshold, indicating that the noise from the unresolved sources will not produce a measurable imprint on the sensitivity curve.
Future developments could include the use of more advanced population synthesis codes and different prescriptions for binary evolution, or the extension of the analysis to larger observational galaxy catalogs and cosmological volumes.
\end{abstract}


\tableofcontents

% Rimuovere se non si vuole la tabella delle figure
% \listoffigures

%\chapter*{Introduction}
\addcontentsline{toc}{chapter}{Introduction}
In the context of the gravitational waves study, there are mainly two possible paths: the first is the \textbf{analysis of existing data} from experiments like Ligo, Virgo and Kagra, with the goal of detecting and characterizing the observable sources within their relative frequency bands;
the second is the \textbf{forecasting approach}, which aims to characterize future experiments in order to better understand what types of sources they could be able to detect and how well.
\vspace{2mm}\\
One of the most important future detectors is the European \textit{Laser Interferometer Space Antenna} (LISA), the first space-based gravitational wave detector. 
LISA will be arranged as an equilateral triangle with 2.5 million kilometers long arms, placed in a heliocentric orbit, and will operate within the frequency range from approximately $0.1\mathrm{mHz}$ to $1\mathrm{Hz}$. 
Among the typical sources that emit gravitational wave signal in this range are the \textbf{compact binaries}, and in our particular case we are interested in \textbf{double white dwarfs} (\textbf{DWDs}).
\vspace{2mm}\\
Many studies have been made to explore the effects of different type of sources on the LISA's sensitivity curve.
For example, we know that the unresolved binary Galactic systems produce a confusion gravitational wave foreground that shapes noticeably the curve close to its lowest threshold.
While the Galactic contribution has been extensively investigated, less attention has been devoted to the possible extragalactic binaries backgrounds.
Therefore, as a natural prosecution, the goal of this work is to estimate the gravitational wave background produced by the \textit{extragalactic} DWDs in the local universe, using the \textbf{COSMIC} code to generate synthetic astrophysical populations with the aid of the galaxies listed in the \textbf{Gravitational Wave Galaxy Catalog} (\textbf{GWGC}), an observational survey, as a reference.
The potential impact of these sources in the tentative to characterize the expected LISA sensitivity and to avoid misinterpreting unresolved signals, are the factors that make this research so important.
\vspace{2mm}
\\  
In \textbf{Theoretical Framework} we introduce the theoretical foundations of gravitational waves, starting from the linearized Einstein field equations, to derive the amplitude of the signal generated by an inspiralling binary system, and define the most important parameters.
We also give a broad theoretical astrophysical framework, briefly discussing the birth and evolution of double white dwarfs, both as single objects, and in binary systems, and the galactic environments that will host them.
We then give a brief overview of how gravitational wave detectors work, trying to better understand what kinds of sources LISA will be able to see, what resolution and what sensitivity it will have and why in particular we are interested in DWDs.
\vspace{2mm}\\
In \textbf{Methods and Data} we introduce the concept of stellar population synthesis and the code COSMIC used for this purpose.
We present its features, main parameters, and pipeline, and explain how it is used to generate full-size astrophysical populations of DWDs.
In order to do this, we introduce the Gravitational Wave Galaxy Gatalog, an observational survay that lists some of the key information we need for the purpose, and explain how we move from there to infer the remaining parameters.
We then explain the procedure to generate the actal fixed populations, scale them to astrophysical full-size galaxies, and compute their gravitational wave signal. 
Moreover, we discuss how to take care of the lack of galaxies due to the catalog's incompleteness and the Milky Way coverage in the Zone Of Avoidance.
Finally, we discuss the LISA frequency resolution problem, and explain how we took it into account.\vspace{2mm}\\
In \textbf{Results and Discussion} we present the results we found with the obtained information and the procedure we described, by plotting the total resulting signal on the LISA's sensitivity curve.
We then discuss what we found, and the possible implications. \vspace{2mm}\\
Finally, in the \textbf{Conclusions and Future Perspectives} we recap the work as a whole, discussing its limitations, assumptions and its possible extensions and follow-ups.

\chapter*{Introduction}
In the context of the gravitational waves study, there are mainly two possible paths: the first is the \textbf{analysis of existing data} from experiments like Ligo, Virgo and Kagra, with the goal of detecting and characterizing the observable sources within their relative frequency bands;
the second approach is the \textbf{forecasting approach}, which aims to characterize future experiments in order to better understand what types of sources they could detect and how well.
\vspace{2mm}\\
One of the most important future detectors is the European \textit{Laser Interferometer Space Antenna} (LISA), the first space-based gravitational wave detector. 
LISA will be arranged as an equilateral triangle with 2.5 million kilometers long arms, placed in a heliocentric orbit, and will operate within the frequency range from approximately $0.1\mathrm{mHz}$ to $1\mathrm{Hz}$. 
Among the typical sources that emit gravitational wave signal in this range are the \textbf{compact binaries}, and in our particular case we are interested in \textbf{double white dwarf binaries} (\textbf{WDBs}).
\vspace{2mm}\\
The goal of this work is to estimate the gravitational wave background produced by the extragalactic WDBs in the local universe, by using the \textbf{COSMIC} code to generate synthetic astrophysical populations to represent the galaxies listed in the \textbf{Gravitational Wave Galaxy Catalog} (\textbf{GWGC}).\vspace{2mm}
\\  
In \textbf{Chapter 1} we will introduce the theoretical foundations of gravitational waves, derive the amplitude of the signal generated by a binary system, and define the most important parameters. Finally, we discuss how to combine the signals from multiple sources, to find a cumulative background.\vspace{2mm}\\
In \textbf{Chapter 2} we will give a brief overview of how gravitational wave detectors work, trying to better understand what kinds of sources LISA will be able to see, what resolution and what sensitivity it will have and why in particular we are interested in WDBs. \vspace{2mm}\\
In \textbf{Chapter 3} we introduce the concept of stellar population synthesis and the code COSMIC used for this purpose. We will introduce its features, main parameters, and pipeline, and explain how it is used to generate full-size astrophysical populations of WDBs.\vspace{2mm}\\ 
In \textbf{Chapter 4} we introduce the GWGC, list the key information it provides and explain how we move from there to infer the remaining parameters that we need.\vspace{2mm}\\
In \textbf{Chapter 5} we use the obtained information to compute the total gravitational wave signal summing the contribution from all the simulated sources, taking into account their spatial distribution, LISA's frequency resolution, and the \textit{zone of avoidance} caused by the milky way.\vspace{2mm}\\
In \textbf{Chapter 6} we plot the total resulting signal on the LISA sensitivity curve, and discuss the results and their possible implications.\vspace{2mm}\\
Finally, in \textbf{Chapter 7} we will draw some conclusions from the work as a whole, discussing its limitations, assumptions and its possible extensions and follow-ups.

\chapter{Gravitational Waves Theory}
After considering, in 1905, the problem of the apparently \textit{instantaneous} propagation of light, with the theory of Special Relativity, in 1916 Albert Einstein considered the problem of the apparently \textit{instantaneous} propagation of gravity through \textit{long distances}, in his theory of General Relativity.
Einstein showed that long-distance interaction arises from the deformation of spacetime caused by massive objects.
Hence, in the "static case", the deviated motion apparently caused by the interaction between two distant masses really is, in fact, a manifestation of spacetime curvature nearby, generated by the presence of the two objects.
The "static case" just depicted, though, treats the curvature as if it had always been there, and doesn't take into account of any variation in the masses, positions or velocities of the two objects, that would induce an evolution to the curvature itself.
In truth, after a change in the mass-energy distribution, the corresponding curvature variation requires its time to reach far distances, and a fascinating prediction of General Relativity is that it propagates in the form of a wave, that travels at the speed of light.



\section{The Flat Spacetime}

\section{Gravitational Waves as Perturbations}
\subsection{Harmonic gauge}
\subsection{The TT gauge}


\section{Motion and geodesics}
\subsection{The motion deviation}
\subsection{The geodesic deviation}

\section{The Quadrupole Approximation}
\subsection{The weak-field, slow-motion approximation}
\subsection{The quadrupole formula}
\subsection{Transform to the TT gauge}

\section{Gravitational waves from a binary system}
\subsection{General solution for circular orbits}
Up to 13.86/7 at page 255 of the book.

\section{Energy carried by a gravitational wave}
\subsection{Stress-energy pseudo-tensor}
\subsection{Gravitational wave luminosity}

\section{Evolution of a compact binary system}
\subsection{Signal from inspiralling compact objects}
Here we get to the actual amplitude we used, and the parameters involved.
\subsection{ASD and multiple sources}

\chapter{LISA}
\section{interferometers}



- Instrument description (what is an interferometer, why in space, how it will be made, orbit)
- frequency band - what will it see?
- frequency resolution
- sensibility curve and ASD meaning
- Why WD choice in particular?

\chapter{COSMIC and Stellar Populations Synthesis}
- COSMIC introduction (reference Breivik, BSE, general information)
- COSMIC pipeline (why it's quick, intro to fixed population idea, match conditions, ...)
- important parameters in COSMIC (the ones user chooses - stress on metallicity relevance)
- population's parameters (the ones COSMIC gives to user)
    - parameter distribution graphs
- Scaling to astrophysical population, what is needed (refer to next chapter

\chapter{GWGC and Galaxy Properties}
As we have seen in \textbf{Chapter 2}, the goal of this work is to simulate the gravitational wave background produced by compact binaries in the local universe.
To do this, we must first generate the populations emitting this signal, and we do this by using COSMIC to replicate the existing, observed galaxies in the vicinity of the Milky Way and simulate their stellar content.
To access the required galactic properties, we rely on the dataset provided in \cite{GWGC}, the \textbf{Gravitational Wave Galaxy Catalog} (\textbf{GWGC}).
This catalog includes a list of $53,255$ galaxies within $100Mpc$ from earth, containing information on sky position, distance, blue magnitude, major and minor diameters, position angle, and galaxy type, currently used for follow-up searches of electromagnetic counterparts from gravitational wave searches.

\section{What it has vs what we need}
In principle, we could generate a separate fixed population for each galaxy in the GWGC and scale it individually.
However, this is simply not practical because of the computational power and time it would require, and therefore we must find a strategy to group them in a few, representative, categories.
As we will show in this section, many of the information in the GWGC can be used to infer the missing astrophysical quantities we need for population synthesis. 
Ultimately, we will find that metallicity is the most suitable parameter for grouping galaxies.
To get there, we follow a chain of empirical relations, starting from the galaxy morphological type, through  a luminosity to mass, and then a mass to metallicity relation. 
This process will enable us to provide COSMIC with the necessary input, found in a consistent and astrophisically motivated way.

\subsection{Galaxy types}

    - T value to Class

\subsection{Mass-luminosity relation}

    - Class to Mass-Luminosity relation -> Mass (Faber \& Gallagher)

\subsection{Mass-metallicity relation}

    - Finding each galaxy's metallicity from the mass (Tremonti - Allende Prieto)


\section{The final fixed populations}

\subsection{Final catalog editing}

    - Graphs of catalog distributions

\subsection{Metallicity categorization}

    - Metallicity bins for GWGC galaxies, and corresponding fixed populations


\chapter{Total Gravitational Wave Signal}
- For each galaxy:
    - Compute right $N_astro$
    - compute each binary's GW signal;
    - bin it to LISA's frequency sensibility bins
    - Plot it on LISA's sensibility curve
- Zone of avoidance: how to consider all the sky


\chapter{Results}
- Plot of spectral distribution of the computed signal
- Analysis of the distribution of the sources
- Eventual implications (none, since it shouldn't be visible)

\chapter{Conclusions and Future Perspectives}
 Recap of the whole Work
- Limits of the Work, assumptions and approximations used
- Possible extensions

\appendix

\chapter{}
\begin{table}
    \caption{In this table are presented all the possible kstar values that the user can choose in COSMIC, and the corresponding astrophysical object and evolutionary state.}\label{tab: kstar values}
    \begin{center}
        \begin{tabular}[]{l|l}
            \toprule
            \multicolumn{1}{c|}{\textbf{kstar value}} & 
            \multicolumn{1}{c}{\textbf{evolutionary state}} \\
            \midrule
            0 & Main Sequence (MS), $< 0.7 M_\odot$ \\
            1 & MS, $>0.7 M_\odot$ \\
            2 & Hertzsprung Gap \\
            3 & First Giant Branch \\
            4 & Core Helium Burning \\
            5 & Early Asymptotic Giant Branch (AGB) \\
            6 & Thermally Pulsing AGB \\
            7 & Naked Helium Star MS \\
            8 & Naked Helium Star Hertzsprung Gap \\
            9 & Naked Helium Star Giant Branch \\
            10 & Helium White Dwarf \\
            11 & Carbon/Oxygen White Dwarf \\
            12 & Oxygen/Neon White Dwarf \\
            13 & Neutron Star \\
            14 & Black Hole \\
            15 & Massless Remnant \\          
            \bottomrule
        \end{tabular}
    \end{center}
\end{table}


% \chapter{}
% Dettagli tecnici sul codice.
% 
% Tabelle di parametri.
% 
% Ulteriori grafici.
% 
% Script di calcolo, se rilevante.


%\chapter{Questa è un'appendice}

%\lipsum[5]

%\section*{Questa è un'altra sezione, ma non viene inserita nell'indice}

%\lipsum[6]

%\bibliography{chapters/Bibliografia}
\nocite{*}
\bibliographystyle{plainnat}
\bibliography{chapters/Bibliografia.bib}

\end{document}
% -----------------------------------------------------------------

